%% ============================================================
\section{Limitations}
\label{sec:limitations}
%% ============================================================

\paragraph{Benchmark coverage.}
We designed CoE and \coeaudit{} to be domain-agnostic, but validating that generality requires evaluation across diverse scientific domains.
Our current experiments focus on systems-optimization tasks (ADRS), where gold-standard evaluators make score verification and specification violation detection straightforward.
Open-ended domains---biology, materials science, theoretical ML---pose harder challenges: evidence chains may involve wet-lab protocols, simulation reproducibility, or proof sketches, each demanding domain-specific verification logic that we have not yet built or tested.
Extending CoE to these settings is a natural next step, though the core abstraction (claims linked to evidence via typed provenance records) should transfer. What changes is the set of integrity checks required.

% \paragraph{Native \cpr{} coverage.}
% Ideally, every claim in a paper would carry machine-verifiable provenance from the moment it is written.
% In practice, \sys{} is the only system that emits structured provenance at claim-production time, making native \cpr{} a \sys{}-only metric.
% Even within \sys{}, our Claim Verifier handles numerical, citation, and methodological claims but does not yet verify conclusion claims reliably---these require reasoning about whether results actually support the broader narrative, a substantially harder problem than checking whether a number matches an evaluator log.
% To enable fair cross-system comparison despite this asymmetry, we rely on integrity checks (\S\ref{sec:coeaudit:checks}), which can be applied uniformly to any system's artifacts in forensic mode, without requiring cooperation from the system under audit.

% \paragraph{\coeaudit{} automation level.}
% We aimed to make \coeaudit{} fully automated, but the current implementation sits on a spectrum.
% Score Verification and Specification Violation Detection are deterministic---they compare hashes and re-run evaluators, with no model in the loop.
% Reference Verification is mostly automated (API lookups against Semantic Scholar), though it uses an LLM disambiguation step when metadata is ambiguous.
% Method-Code Alignment, however, is fundamentally LLM-judged: it asks whether a natural-language method description matches submitted code, a task that resists deterministic automation.
% We report the models and prompts used (\Cref{tab:audit_procedure}), but we have not yet measured inter-rater reliability across different judge models---an important step before \coeaudit{} scores can be treated as ground truth rather than informed estimates.

\paragraph{Reference verification depth.}
Our Reference Verification checks whether cited references \emph{exist}---a necessary condition that already catches a surprising number of failures (\S\ref{sec:exp:failures}, Case~2).
However, existence is far from sufficient: a real citation can still be used to support a claim the cited paper never made.
Full reference verification would require passage-level natural language inference against the cited paper's text, effectively asking ``does this source actually say what the citing paper claims it says?''
This is a known open problem in scholarly NLI, and we leave it to future work, noting that even existence-level checking already reveals meaningful architectural differences between systems.

% Hidden for internal review -- restore for submission if needed:
\paragraph{Automated review as proxy.}
ScholarPeer serves as a scalable proxy for review quality but does not replace human expert evaluation.
LLM reviewers are systematically blind to certain failure modes (e.g., domain-specific score interpretation, specification violation detection).
\coeaudit{} addresses some of these blind spots but is itself limited to structural integrity, not scientific novelty or significance.

% \paragraph{PI $\to$ solver causal link.}
% The Problem Investigator provides grounded literature context, but its effect on solver performance is mixed.
% We do not claim PI improves solver performance---only that it provides grounded context for paper writing and reference integrity.

\paragraph{Fairness of baseline comparison.}
We adapted four open-source systems to ADRS under conditions as uniform as possible---identical model backbone, matching wall-clock limits, and three seeds per task (Appendix~\ref{app:baseline_adaptation}).
%Evaluator-call budgets are comparable: \sys{}'s base configuration makes ${\sim}100$ evaluator calls (25 tree nodes $\times$ ${\sim}4$ evaluated solver versions each), matching AdaEvolve's $T\!=\!100$ iteration budget~\citep{cemri2026adaevolve}; our adapted baselines receive up to 20 solver iterations per task.
%What differs is search structure---\sys{} distributes evaluations across a tree with ideation and pruning between iterations, adding LLM overhead (ideator, auditor, ablation, checkpointer) not counted in the evaluator budget---so the total model inference per system varies even when evaluator-call counts match.
No third-party system was designed for ADRS, and adaptation inevitably involves judgment calls.
We erred on the side of generosity (e.g., giving ARC $6.7\times$ its default iteration budget, re-running infrastructure crashes but never re-running to improve scores), yet we cannot rule out that a system's original authors would achieve better results with deeper tuning.
Cross-system comparisons should be read as ``given a good-faith, equal-resource adaptation'' rather than ``definitive system ranking.''

\paragraph{Audit false negatives.}
All I1--I3 flagged positives were manually verified by human reviewers, ensuring that the integrity failures reported in Table~\ref{tab:integrity} are real (no false positives).
However, we did not systematically bound false negatives: integrity failures that our checks fail to detect certainly exist, and the true failure rate across all systems is likely higher than reported.
For I4 (Method-Code Alignment), human reviewers validated a sample of the LLM judgments, but the reported results reflect majority-vote scores without systematic correction.
This is an inherent limitation of any finite audit protocol, not specific to \coeaudit{}.

\paragraph{Benchmark scope and depth.}
ADRS tasks reduce systems research problems to single-metric optimization---submit a solver, receive a score.
Real systems papers involve problem formulation, workload characterization, multi-dataset analysis, and deployment tradeoffs that our pipeline does not attempt.
As a result, ``competitive solver performance on ADRS'' should not be equated with ``competitive systems research.''
Extending to multi-benchmark synthesis and deeper experimental analysis---where the system reasons about \emph{why} a solution works, not just \emph{whether} it scores well---is a natural next direction.

\paragraph{Broader impacts.}
Autonomous research systems that produce full papers create both opportunities and risks.
On the positive side, structured provenance (CoE) and systematic auditing (\coeaudit{}) make integrity failures detectable at a scale that manual review cannot match---every number, citation, and method claim can be traced to its source artifact.
On the negative side, the same capability lowers the barrier to generating plausible-looking scientific papers at volume, potentially flooding review pipelines or producing results that appear rigorous but contain subtle errors outside the audit's scope.
Our integrity checks mitigate this by making verifiable papers distinguishable from unverifiable ones, but they cover structural integrity, not scientific correctness or novelty.
We believe that transparency tools like CoE Audit should be developed alongside generation capabilities, so that the research community can verify AI-generated claims rather than taking them on trust.

% \paragraph{Forensic CPR extraction noise.}
% Forensic claim extraction from \LaTeX{} is inherently lossy.
% For \sys{}, the gap between native and forensic CPR quantifies this noise.
% Forensic CPR comparisons across systems are fair (all systems face the same extraction noise), but absolute values underestimate true verifiability.
